\chapter{Data, Observations, and Methods}

This chapter demonstrates key \LaTeX{} formatting features available in
this template: running text, mathematical expressions, footnotes, a
table, and a figure placeholder.

\section{Sample and Survey Overview}

Describe your dataset here \citep{example:2021}. Inline citations use
\citet{example:2020}; parenthetical citations use
\citep[e.g.,][]{example:2020}.%
\footnote{Footnotes look like this. Keep them concise.}

A summary of the sample properties is given in Table~\ref{tab:sample}.
All results assume a flat $\Lambda$CDM cosmology with
$H_0 = 70\,\mathrm{km\,s^{-1}\,Mpc^{-1}}$, $\Omega_\mathrm{m} = 0.3$.


\subsection{Selection Criteria}

Objects were selected according to the following criterion:
\begin{equation}
    \log \left( \frac{L}{L_\odot} \right) > 10 ,
    \label{eq:selection}
\end{equation}
yielding a final sample of $N$ sources spanning a redshift range
$z_\mathrm{min} < z < z_\mathrm{max}$.


\section{Analysis Method}

Each source was modelled using a S\'{e}rsic profile
\citep{example:2020},
\begin{equation}
    I(r) = I_e \exp\left\{ -b_n \left[
        \left(\frac{r}{r_e}\right)^{1/n} - 1 \right] \right\} ,
    \label{eq:sersic}
\end{equation}
where $r_e$ is the half-light radius, $I_e$ is the surface brightness
at $r_e$, and $n$ is the S\'{e}rsic index.

Figure~\ref{fig:ch2_example} shows an illustrative example.

\begin{figure}[htbp]
    \centering
    % \includegraphics[width=0.9\textwidth]{your_figure.pdf}
    \framebox[0.9\textwidth]{\rule{0pt}{5cm}%
        \parbox{0.85\textwidth}{\centering\small
        Figure placeholder --- replace with
        \texttt{\textbackslash includegraphics\{your\_figure.pdf\}}}}
    \caption[Short caption for list of figures]{Full figure caption goes
        here. Describe all panels. Left: \ldots{} Right: \ldots}
    \label{fig:ch2_example}
\end{figure}


\begin{table}[htbp]
\centering
\caption[Sample summary]{Summary of the sample used in this work.}
\label{tab:sample}
\begin{tabular}{lcc}
\toprule
Property & Value & Unit \\
\midrule
Total sources   & $N$           & ---       \\
Redshift range  & $z_1$--$z_2$  & ---       \\
Median redshift & $\bar{z}$     & ---       \\
Median $\log L$ & $X.X$         & $L_\odot$ \\
\bottomrule
\end{tabular}
\end{table}


\section{Summary}

\begin{enumerate}
    \item The sample contains $N$ sources selected via
          Equation~(\ref{eq:selection}).
    \item Surface-brightness profiles were fitted with the S\'{e}rsic
          model (Equation~\ref{eq:sersic}).
    \item Results are presented in Chapter~3.
\end{enumerate}
